\section{The genuine common convex profiles have the common leading functions}

Let $k$ be a field and let $f:Y\to\operatorname{Spec}k[t]$ be the
original proper dominant morphism with integral $Y$. Let
$i:C\hookrightarrow Y$ be the original integral closed curve, assume
$C$ locally Noetherian and $\dim C\le1$, and put $\eta=i(\eta_C)$.
One actual Noetherian point with a genuine prime component equation is
given, as in the existing common-profile producer. Retain the actual
verticality condition $t\in(f\eta)$.

Write $D=\mathcal O_{Y,\eta}$, $K=Y.\mathrm{functionField}$,
$F=C.\mathrm{functionField}$, and let $t_\eta\in D$ and
$t_x\in S_x=\mathcal O_{Y,i(x)}$ be the ORIGINAL germs of the original
base parameter. Every localization, weighted DVR, fraction embedding,
residue map, Gauss point and native function below is the SAME old one.

\begin{lemma}[The exact original field and parameter diagrams]
\label{lem:gauss-common-leading-original-field-parameter}
\uses{lem:gauss-native-original-stalk-point-sequence}
\lean{EnvelopeProperty.GaussCommonLeadingProfiles.genericFieldMap_algebraMap,EnvelopeProperty.GaussCommonLeadingProfiles.genericFieldMap_parameter}
\leanok
For an original $p\in Y$ let $J_p:\operatorname{Frac}\mathcal O_{Y,p}\simeq K$
be the original canonical fraction-field isomorphism. The field map
used by the actual old profiles is EXACTLY $J_p^{-1}$. For EVERY original
$s\in\mathcal O_{Y,p}$ it sends the original image of $s$ in $K$ to the
original fraction coefficient $s/1$.
For ANY original points $p,q\in Y$, their actual base-parameter germs
therefore satisfy
\[
 J_p^{-1}(\iota_q(t_q))=\iota_p(t_p).
\]
In particular the parameter fraction used by the original weighted
profiles is the SAME fraction used by the common-leading producer at
$q=\eta$. No field-matching premise is supplied.
\end{lemma}
\begin{proof}
The old profile construction defines its field map by the inverse of the
canonical fraction-ring isomorphism for the SAME original stalk. The
canonical isomorphism commutes with every original stalk coefficient;
applying its inverse proves the all-coefficient identity. Alternatively,
these coefficient identities uniquely determine its whole-field map,
by the fraction-field universal property.

The base parameter is the actual global section obtained by applying
the original $f$ to $t$. The actual germ-to-function-field square identifies
the image of its germ at $q$ with this SAME section's original generic
value. The old profile germ-square then takes this SAME value to the
original fraction coefficient of its germ at $p$. Combining the two
actual squares proves the displayed identity for every $p,q$.
\end{proof}

\begin{lemma}[The actual unnormalized weighted parameter order]
\label{lem:gauss-common-leading-raw-parameter-order}
\uses{def:gauss-surface-same-field-native-orders,lem:gauss-surface-parameter-positive-vertical-order,lem:gauss-surface-original-vertical-fraction-identities,thm:gauss-regular-surface-original-dvr}
\lean{EnvelopeProperty.GaussCommonLeadingProfiles.weightedOrder_parameter}
\leanok
Fix genuine original local parameters $\pi,z$ of a Noetherian local
domain $S$, with $\pi$ prime, $\pi\nmid z$ and
$\mathfrak m_S=(\pi,z)$. For an original nonzero $t_S\in S$ divisible
by $\pi$, put
$d_N=\operatorname{weightedParameterOrder}_N(t_S)$.
For EVERY positive original $N$,
\[
 \operatorname{weightedOrder}_N(\iota_S(t_S))=\exp(d_N).
\]
This is an identity of the original normalized INTEGER DVR order,
before division by the parameter order or passage to any real norm.
\end{lemma}
\begin{proof}
Let $V_N$ be the existing actual weighted DVR and keep its actual original
fraction inclusion in $\operatorname{Frac}S$. The existing all-coefficient
inclusion formula identifies $\iota_S(t_S)$ with the image of the actual
coefficient of $t_S$ in $V_N$.
The actual weighted-parameter specification proves that the additive
DVR order of this coefficient is exactly the finite natural number
$d_N$. The native whole-field order is the inverse of the valuation at
the SAME actual maximal ideal. On an original DVR coefficient that
valuation is the inverse integer exponential of its actual additive
order. Substituting $d_N$ and taking the inverse yields $\exp(d_N)$.
All scalar actions are the literal original coefficient actions of the
existing weighted-order definition. No scale or norm comparison enters.
\end{proof}

\begin{lemma}[An original leading tuple determines its actual raw affine tail]
\label{lem:gauss-common-leading-actual-tuple-tail}
\uses{thm:gauss-surface-original-fraction-eventual-orders,def:gauss-surface-same-field-native-orders,lem:gauss-surface-original-vertical-fraction-identities,lem:gauss-surface-original-residual-vertical-dvrs}
\lean{EnvelopeProperty.GaussCommonLeadingProfiles.eventual_weightedOrder_of_leading}
\leanok
For the SAME original local geometry let $g\ne0$ lie in
$\operatorname{Frac}S$. Suppose an actual original $u\in S_{(\pi)}^\times$
and integers $a,b$ have the genuine original properties
\[
 g=\iota_T(u)\pi^a,\qquad
 \operatorname{residualOrder}(\operatorname{res}u)=\exp(b).
\]
Then for all sufficiently large positive $N$ the UNCHANGED actual native
order satisfies
\[
 \operatorname{weightedOrder}_N(g)=\exp(aN+b).
\]
The final producer below obtains these actual intermediate data internally
from the old producers, rather than requesting them as original inputs.
\end{lemma}
\begin{proof}
Apply the old all-fraction theorem to the SAME $g$. It produces
$a',b',u'$ with its actual factorization, actual vertical and residual
orders, and the actual eventual weighted-order tail
$\exp(a'N+b')$.
Apply the native original vertical order to the stated factorization of
$g$. An actual vertical unit has order one, and the original $\pi/1$
is a true uniformizer. Thus this factorization gives order $\exp(a)$.
Comparing with the old theorem's $\exp(a')$ and using injectivity of the
integer exponential gives $a'=a$.

Cancel the same nonzero power of the original $\pi$ in the two genuine
whole-field factorizations. The original vertical fraction inclusion is
injective, so the actual ring elements $u'$ and $u$ coincide. Their unit
group elements consequently coincide as well. Their canonical residues
and native residual orders are now literally the same. Comparing the
stated $\exp(b)$ with the old returned $\exp(b')$ gives $b'=b$.
Substitute these TWO exact integer equalities into the old actual
eventual tail. This derives the asserted raw law from the actual
unchanged maps and units.
\end{proof}

\begin{lemma}[Two adjacent raw orders determine both affine coefficients]
\label{lem:gauss-common-leading-affine-uniqueness}
\lean{EnvelopeProperty.GaussCommonLeadingProfiles.affine_coefficients_eq_of_eventually}
\leanok
Let $a,b,A,B\in\mathbf Z$. If eventually in the original natural index $m$
\[
 \exp(a(m+1)+b)=\exp(A(m+1)+B),
\]
then $a=A$ and $b=B$.
\end{lemma}
\begin{proof}
Choose $M$ beyond the actual eventual threshold. Apply the equality at
the adjacent original indices $M$ and $M+1$. Injectivity of the integer
exponential gives
\[
 a(M+1)+b=A(M+1)+B,\qquad
 a(M+2)+b=A(M+2)+B.
\]
Subtracting the first equality from the second gives $a=A$. Substitution
in the first and cancellation of the common integer summand gives
$b=B$. This uses the raw integer orders with the SAME original index;
an equality of normalized ratios alone would not imply these conclusions.
\end{proof}

\begin{theorem}[The actual old profile coefficients are the original leading coefficients]
\label{thm:gauss-common-leading-profile-coefficients}
\uses{lem:gauss-common-leading-original-field-parameter,lem:gauss-common-leading-raw-parameter-order,lem:gauss-common-leading-actual-tuple-tail,lem:gauss-common-leading-affine-uniqueness}
\lean{EnvelopeProperty.GaussCommonLeadingProfiles.profile_orders_eq,EnvelopeProperty.GaussCommonLeadingProfiles.parameter_profile_orders_eq}
\leanok
Retain an actual original regular-unit family and the genuine outputs of
the old original profile producer. Denote its local data by
$e_{\rm old},c_{\rm old}\in\mathbf N$, $A_\ell,B_\ell\in\mathbf Z$,
and its actual original vertical units by $v_\ell$.
Apply the common-leading producer to the SAME original parameter
$t_\eta$ and SAME global rational functions represented by these units.
Denote its actual data at $x$ by $e>0$, $a_\ell,b_{\ell,x},c_x$.
Then the SAME actual outputs satisfy
\[
 A_\ell=a_\ell,\qquad B_\ell=b_{\ell,x},\qquad
 e_{\rm old}=e,\qquad (c_{\rm old}:\mathbf Z)=c_x.
\]
These are equalities of the actual integer coefficients, including the
constant of the original parameter law, before any real normalization.
\end{theorem}
\begin{proof}
The first lemma identifies the old and new field elements using the
actual original coefficient squares. For each original regular unit,
both original producers give genuine local fraction factorizations and
genuine native residual orders. Apply the actual-tuple tail lemma to
each tuple. Their left sides are the SAME unnormalized original
$\operatorname{weightedOrder}_N$ of the SAME original fraction.
Intersect the two eventual tails and apply adjacent-index affine
uniqueness. It gives both $A_\ell=a_\ell$ and $B_\ell=b_{\ell,x}$.

For the parameter, the old actual profile gives, eventually,
$d_{m+1}=e_{\rm old}(m+1)+c_{\rm old}$.
The actual-tuple tail applied to the new parameter tuple gives the
actual native order $\exp(e(m+1)+c_x)$.
The exact raw parameter-order lemma identifies the latter left side
with $\exp(d_{m+1})$, for the SAME coefficient $t_x$ and SAME actual
weighted DVR. Thus the two actual raw affine laws have equal integer
exponentials eventually. Adjacent-index uniqueness gives the claimed
$e_{\rm old}=e$ and $(c_{\rm old}:\mathbf Z)=c_x$, with no common
rescaling left undetermined.

The two comparisons are independent, so the parameter statement remains
valid when there are no regular-unit coordinates. Neither comparison
uses a normalized real value or assumes equality of leading functions.
\end{proof}

\begin{theorem}[The original common active profile uses one family of global leading functions]
\label{thm:gauss-common-global-leading-convex-profiles}
\uses{thm:gauss-common-original-convex-family,thm:gauss-common-leading-profile-coefficients,lem:gauss-common-leading-original-field-parameter,thm:gauss-leading-corrected-curve-order,alggeom:original-stalk-coefficient-map,alggeom:original-stalk-coefficient-injective,alggeom:original-stalk-coefficient-center,analysis:actual-finite-affine-active-max,scheme:native-regular-unit-local-psh}
\lean{EnvelopeProperty.GaussCommonLeadingProfiles.exists_common_globalLeading_profiles}
\leanok
Let $u$ satisfy the original $\operatorname{IsConvexPLOn}(U,u)$ on the
original native space at the paper Gauss radius. The SAME original
common-point producer gives $q_C$ with true center $\eta$ and original
generic support. Under the honest original applicability condition
$q_C\in U$, choose ONCE the actual $W,V,a_1,\ldots,a_n,I$ and the
nonempty active set $I_0$ of the original common-profile producer.
Internally produce ONE $e>0$, the SAME global vertical exponents
$a_1',\ldots,a_n'\in\mathbf Z$ and nonzero actual functions
$B_1,\ldots,B_n,B_t\in F$, ALL before choosing $x\in C$.

For EVERY nongeneric original $x$ with the genuine original local
component geometry define
\[
 \gamma_{\ell,x}=\operatorname{ord}_x(B_\ell)
        -\frac{a_\ell'}e\operatorname{ord}_x(B_t),\qquad
 s_{j,x}=-\sum_\ell j_\ell\gamma_{\ell,x},\qquad
 D_x=\max_{j\in I_0}s_{j,x}.
\]
The original actual weighted points $q_{x,m}$ and scales
$\varepsilon_{x,m}$ retain their original definitions. Eventually,
\[
 u(q_{x,m})=u(q_C)+D_x\varepsilon_{x,m},\qquad
 \frac{u(q_{x,m})^{\rm real}-u(q_C)^{\rm real}}
      {\varepsilon_{x,m}}=D_x,
\]
and the actual quotient tends to this SAME $D_x$.
All original source-family values, indices, active indices and global
leading functions are fixed before varying $x$. No profile-coefficient,
residue, field, degree or norm comparison is an original input.
\end{theorem}
\begin{proof}
Use the old common-profile producer once. It chooses the actual $q_C$.
Under its original applicability condition, it applies the original
source convexity hypothesis once and supplies the SAME $W,V,a,I,I_0$
with all its actual pointwise values and local eventual profiles.

The actual old coefficient map is the original section map followed by
the genuine germ map. Its actual contraction of the generic-stalk
maximal ideal is the original base prime. Original verticality therefore
puts $t_\eta$ in the true maximal ideal. Original proper dominance gives
injectivity of that SAME coefficient map, so $t_\eta\ne0$.
Apply the frozen actual common-leading producer to this SAME actual
parameter. It internally supplies the uniformizer, $e>0$, and the
actual parameter leading function. Each original unit $a_\ell$ gives
its SAME nonzero global rational function by the original generic
section map. Choose its global factor and leading function using that
same producer. These finite choices precede all original local points.
Also retain its parameter tuple by applying its all-function statement
to the actual nonzero constant function $1$. Thus the parameter comparison
is available even when the original regular-unit family is empty; no
new source presentation or active index is selected.

Now fix an original $x$ and its actual local parameters. Retain the old
common-profile producer's true point equality, actual unit tuple and
actual eventual parameter law. Apply the preceding coefficient theorem
to those SAME old tuples and the newly produced original leading
tuples. This identifies their actual integer coefficients, including
the actual $c_{\rm old}$ in this specific profile. The unchanged old
local correction is therefore exactly $\gamma_{\ell,x}$ above.
Finite summation identifies every old original affine slope with
$s_{j,x}$. The intercepts and active set were fixed before $x$, so the
actual old active maximum is precisely the stated $D_x$.
Substitute this equality into the old actual function profile, exact
actual quotient and its true limit. Both old function values are
finite, so their real-part identities are unchanged. This proves every
assertion for the SAME original points, scales and function.
\end{proof}

No principal-degree sum, residual nefness, global SNC existence,
LocalPositivity or final envelope assembly is asserted here. The
purpose is the genuine connection of the already produced original
convex profiles to actual whole-curve leading functions, with exact
unnormalized coefficient control.
